\subsection{Active}\label{active}
\begin{figure}[H]
	\centering
	\begin{tikzpicture}[node distance = 3cm, auto, thick,scale=\CrossAndTopSectionBig, every node/.style={transform shape}]
		% substrate
\fill[substrate] (0,0) rectangle (20,2);
\node at (2,0.5) {Silicon substrate};
%trenches
\fill[isolationoxide] (0,0.75) rectangle (1,2);
\fill[isolationoxide] (8.5,0.75) rectangle (11.5,2);
\fill[isolationoxide] (19,0.75) rectangle (20,2);
% n-well
\fill[nwell] (1,0.75) rectangle (8.5,2);
\node at (5.75,1) {N-Well};
% gate oxide
\fill[LightGray] (4.8,2) rectangle (6.7,2.3);
\fill[LightGray] (13.3,2) rectangle (15.2,2.3);
% gate poly
\fill[BrickRed] (4.8,2.3) rectangle (6.7,2.6);
\fill[BrickRed] (13.3,2.3) rectangle (15.2,2.6);
	\end{tikzpicture}
	\begin{tikzpicture}[node distance = 3cm, auto, thick,scale=\CrossAndTopSectionBig, every node/.style={transform shape}]
		\fill[YellowOrange] (0,0) rectangle (20,10);

% n-well
\fill[Goldenrod] (1,1.25) rectangle (8.5,7.5);

% field oxide
\fill[DarkGray] (0,0) rectangle (1,12);
\fill[DarkGray] (8.5,0) rectangle (11.5,12);
\fill[DarkGray] (19,0) rectangle (20,12);
\fill[DarkGray] (0,0) rectangle (20,1.25);
\fill[DarkGray] (0,7.5) rectangle (20,12);
% poly
\fill[BrickRed] (4.8,1.75) rectangle (6.7,9);
\fill[BrickRed] (13.3,1.75) rectangle (15.2,9);
\fill[BrickRed] (6.7,8) rectangle (15.2,9);
	\end{tikzpicture}
	\caption{Active area geometry target}
\end{figure}

\subsubsection{Gate oxide growth}
The thickness of the gate oxide depends the capacity (slew rate) of the transistor.
The thinner the layer is, the steeper the edges of the CMOS circuitry will be, however also the threshold voltage will be reduced the thinner the gate oxide gets.

\begin{figure}[H]
	\centering
	\begin{tikzpicture}[node distance = 3cm, auto, thick,scale=\CrossSectionOnly, every node/.style={transform shape}]
		% substrate
\fill[substrate] (0,0) rectangle (20,2);
\node at (2,0.5) {Silicon substrate};
%trenches
\fill[isolationoxide] (0,0.75) rectangle (1,2);
\fill[isolationoxide] (8.5,0.75) rectangle (11.5,2);
\fill[isolationoxide] (19,0.75) rectangle (20,2);
% n-well
\fill[nwell] (1,0.75) rectangle (8.5,2);
\node at (5.75,1) {N-Well};
	\end{tikzpicture} \\
	\includegraphics[scale=0.01]{down_arrow.png} \\
	\begin{tikzpicture}[node distance = 3cm, auto, thick,scale=\CrossSectionOnly, every node/.style={transform shape}]
		% substrate
\fill[YellowOrange] (0,0) rectangle (20,2);
\node at (2,0.5) {Si (p-type)};
% n-well
\fill[Goldenrod] (1.25,0.75) rectangle (8.25,2);
\node at (4.75,1) {N-Well};
% gate oxide layer
\fill[LightGray] (0,2) rectangle (20,2.1);
%field oxides:
\fill[DarkGray] (0,2) rectangle (1,4);
\fill[DarkGray] (8.5,2) rectangle (11.5,4);
\fill[DarkGray] (19,2) rectangle (20,4);
% channel stop
\fill[RedOrange] (0,1.5) rectangle (1,2);
\fill[RedOrange] (8.5,1.5) rectangle (11.5,2);
\fill[RedOrange] (19,1.5) rectangle (20,2);
	\end{tikzpicture}
	\caption{Gate oxide layer growth}
\end{figure}

\subsubsection{Polysilicon growth}
In this step we need to grow the material for making the gate contact. We use polysilicon instead of metal because polysilicon doesn't melt away during annealing and other high temperature steps.
\begin{figure}[H]
	\centering
	\begin{tikzpicture}[node distance = 3cm, auto, thick,scale=\CrossSectionOnly, every node/.style={transform shape}]
		% substrate
\fill[substrate] (0,0) rectangle (20,2);
\node at (2,0.5) {Silicon substrate};
%trenches
\fill[isolationoxide] (0,0.75) rectangle (1,2);
\fill[isolationoxide] (8.5,0.75) rectangle (11.5,2);
\fill[isolationoxide] (19,0.75) rectangle (20,2);
% n-well
\fill[nwell] (1,0.75) rectangle (8.5,2);
\node at (5.75,1) {N-Well};
\fill[LightGray] (0,2) rectangle (20,2.3);
	\end{tikzpicture} \\
	\includegraphics[scale=0.01]{down_arrow.png} \\
	\begin{tikzpicture}[node distance = 3cm, auto, thick,scale=\CrossSectionOnly, every node/.style={transform shape}]
		% substrate
\fill[substrate] (0,0) rectangle (20,2);
\node at (2,0.5) {Silicon substrate};
%trenches
\fill[isolationoxide] (0,0.75) rectangle (1,2);
\fill[isolationoxide] (8.5,0.75) rectangle (11.5,2);
\fill[isolationoxide] (19,0.75) rectangle (20,2);
% n-well
\fill[nwell] (1,0.75) rectangle (8.5,2);
\node at (5.75,1) {N-Well};
\fill[LightGray] (0,2) rectangle (20,2.3);
\fill[BrickRed] (0,2.3) rectangle (20,2.6);
	\end{tikzpicture}
	\caption{Polysilicon layer growth}
\end{figure}
The LPCVD polysilicon is deposited at 620\degree C and then doped using $POCl_3$  with 1100 ppm partials pressure at 950\degree C for 30 minutes.

\subsubsection{Patterning}
\begin{figure}[H]
	\centering
	\begin{tikzpicture}[node distance = 3cm, auto, thick,scale=\CrossAndTopSection, every node/.style={transform shape}]
		% substrate
\fill[substrate] (0,0) rectangle (20,2);
\node at (2,0.5) {Silicon substrate};
%trenches
\fill[isolationoxide] (0,0.75) rectangle (1,2);
\fill[isolationoxide] (8.5,0.75) rectangle (11.5,2);
\fill[isolationoxide] (19,0.75) rectangle (20,2);
% n-well
\fill[nwell] (1,0.75) rectangle (8.5,2);
\node at (5.75,1) {N-Well};

% gate oxide layer
\fill[LightGray] (0,2) rectangle (20,2.3);
% poly layer
\fill[BrickRed] (0,2.3) rectangle (20,2.6);

	\end{tikzpicture}
	\begin{tikzpicture}[node distance = 3cm, auto, thick,scale=\CrossAndTopSection, every node/.style={transform shape}]
		\fill[BrickRed] (0,0) rectangle (20,10);
	\end{tikzpicture} \\
	\includegraphics[scale=0.01]{down_arrow.png} \\
	\begin{tikzpicture}[node distance = 3cm, auto, thick,scale=\CrossAndTopSection, every node/.style={transform shape}]
		% substrate
\fill[substrate] (0,0) rectangle (20,2);
\node at (2,0.5) {Silicon substrate};
%trenches
\fill[isolationoxide] (0,0.75) rectangle (1,2);
\fill[isolationoxide] (8.5,0.75) rectangle (11.5,2);
\fill[isolationoxide] (19,0.75) rectangle (20,2);
% n-well
\fill[nwell] (1,0.75) rectangle (8.5,2);
\node at (5.75,1) {N-Well};

% gate oxide layer
\fill[LightGray] (0,2) rectangle (20,2.3);
% poly layer
\fill[BrickRed] (0,2.3) rectangle (20,2.6);

% resist
\fill[orange] (4.8,2.6) rectangle (6.7,3.2);
\fill[orange] (13.3,2.6) rectangle (15.2,3.2);

	\end{tikzpicture}
	\begin{tikzpicture}[node distance = 3cm, auto, thick,scale=\CrossAndTopSection, every node/.style={transform shape}]
		\fill[BrickRed] (0,0) rectangle (20,10);
% resist
\fill[orange] (4.8,1.75) rectangle (6.7,9);
\fill[orange] (13.3,1.75) rectangle (15.2,9);
\fill[orange] (6.7,8) rectangle (15.2,9);
	\end{tikzpicture}
	\caption{Polysilicon contact mask}
\end{figure}

\subsubsection{Etching}
Because the exact shape of the gate contact is required for a reproducible property characterization of the transistor geometry, dry etching is being used for etching the poly-oxide layer stack.
\begin{figure}[H]
	\centering
	\begin{tikzpicture}[node distance = 3cm, auto, thick,scale=\CrossAndTopSection, every node/.style={transform shape}]
		% substrate
\fill[YellowOrange] (0,0) rectangle (20,2);
\node at (2,0.5) {Si (p-type)};
% n-well
\fill[Goldenrod] (1.25,0.75) rectangle (8.25,2);
\node at (4.75,1) {N-Well};
% gate oxide layer
\fill[LightGray] (0,2) rectangle (20,2.1);
% poly layer
\fill[BrickRed] (0,2.1) rectangle (20,2.2);
% resist
\fill[orange] (4.8,2.2) rectangle (6.7,2.8);
\fill[orange] (13.3,2.2) rectangle (15.2,2.8);
%field oxides:
\fill[DarkGray] (0,2) rectangle (1,4);
\fill[DarkGray] (8.5,2) rectangle (11.5,4);
\fill[DarkGray] (19,2) rectangle (20,4);
% channel stop
\fill[RedOrange] (0,1.5) rectangle (1,2);
\fill[RedOrange] (8.5,1.5) rectangle (11.5,2);
\fill[RedOrange] (19,1.5) rectangle (20,2);
	\end{tikzpicture}
	\begin{tikzpicture}[node distance = 3cm, auto, thick,scale=\CrossAndTopSection, every node/.style={transform shape}]
		\fill[BrickRed] (0,0) rectangle (20,10);
% resist
\fill[orange] (4.8,1.75) rectangle (6.7,9);
\fill[orange] (13.3,1.75) rectangle (15.2,9);
\fill[orange] (4.8,8) rectangle (15.2,9);
	\end{tikzpicture} \\
	\includegraphics[scale=0.01]{down_arrow.png} \\
	\begin{tikzpicture}[node distance = 3cm, auto, thick,scale=\CrossAndTopSection, every node/.style={transform shape}]
		% substrate
\fill[substrate] (0,0) rectangle (20,2);
\node at (2,0.5) {Silicon substrate};
%trenches
\fill[isolationoxide] (0,0.75) rectangle (1,2);
\fill[isolationoxide] (8.5,0.75) rectangle (11.5,2);
\fill[isolationoxide] (19,0.75) rectangle (20,2);
% n-well
\fill[nwell] (1,0.75) rectangle (8.5,2);
\node at (5.75,1) {N-Well};

% gate oxide
\fill[LightGray] (4.8,2) rectangle (6.7,2.3);
\fill[LightGray] (13.3,2) rectangle (15.2,2.3);
% gate poly
\fill[BrickRed] (4.8,2.3) rectangle (6.7,2.6);
\fill[BrickRed] (13.3,2.3) rectangle (15.2,2.6);

% resist
\fill[orange] (4.8,2.6) rectangle (6.7,3.2);
\fill[orange] (13.3,2.6) rectangle (15.2,3.2);
	\end{tikzpicture}
	\begin{tikzpicture}[node distance = 3cm, auto, thick,scale=\CrossAndTopSection, every node/.style={transform shape}]
		\fill[YellowOrange] (0,0) rectangle (20,10);

% n-well
\fill[Goldenrod] (1,1.25) rectangle (8.5,7.5);

% field oxide
\fill[DarkGray] (0,0) rectangle (1,12);
\fill[DarkGray] (8.5,0) rectangle (11.5,12);
\fill[DarkGray] (19,0) rectangle (20,12);
\fill[DarkGray] (0,0) rectangle (20,1.25);
\fill[DarkGray] (0,7.5) rectangle (20,12);
% resist
\fill[orange] (4.8,1.75) rectangle (6.7,9);
\fill[orange] (13.3,1.75) rectangle (15.2,9);
\fill[orange] (4.8,8) rectangle (15.2,9);
	\end{tikzpicture}
	\caption{Polysilicon contact etched}
\end{figure}
$S F_6$ and $C_2 Cl F_5$

\subsubsection{Cleaning}
\begin{figure}[H]
	\centering
	\begin{tikzpicture}[node distance = 3cm, auto, thick,scale=\CrossAndTopSection, every node/.style={transform shape}]
		% substrate
\fill[substrate] (0,0) rectangle (20,2);
\node at (2,0.5) {Silicon substrate};
%trenches
\fill[isolationoxide] (0,0.75) rectangle (1,2);
\fill[isolationoxide] (8.5,0.75) rectangle (11.5,2);
\fill[isolationoxide] (19,0.75) rectangle (20,2);
% n-well
\fill[nwell] (1,0.75) rectangle (8.5,2);
\node at (5.75,1) {N-Well};

% gate oxide
\fill[LightGray] (4.8,2) rectangle (6.7,2.3);
\fill[LightGray] (13.3,2) rectangle (15.2,2.3);
% gate poly
\fill[BrickRed] (4.8,2.3) rectangle (6.7,2.6);
\fill[BrickRed] (13.3,2.3) rectangle (15.2,2.6);

% resist
\fill[orange] (4.8,2.6) rectangle (6.7,3.2);
\fill[orange] (13.3,2.6) rectangle (15.2,3.2);
	\end{tikzpicture}
	\begin{tikzpicture}[node distance = 3cm, auto, thick,scale=\CrossAndTopSection, every node/.style={transform shape}]
		\fill[YellowOrange] (0,0) rectangle (20,10);

% n-well
\fill[Goldenrod] (1,1.25) rectangle (8.5,7.5);

% field oxide
\fill[DarkGray] (0,0) rectangle (1,12);
\fill[DarkGray] (8.5,0) rectangle (11.5,12);
\fill[DarkGray] (19,0) rectangle (20,12);
\fill[DarkGray] (0,0) rectangle (20,1.25);
\fill[DarkGray] (0,7.5) rectangle (20,12);
% resist
\fill[orange] (4.8,1.75) rectangle (6.7,9);
\fill[orange] (13.3,1.75) rectangle (15.2,9);
\fill[orange] (4.8,8) rectangle (15.2,9);
	\end{tikzpicture} \\
	\includegraphics[scale=0.01]{down_arrow.png} \\
	\begin{tikzpicture}[node distance = 3cm, auto, thick,scale=\CrossAndTopSection, every node/.style={transform shape}]
		% substrate
\fill[substrate] (0,0) rectangle (20,2);
\node at (2,0.5) {Silicon substrate};
%trenches
\fill[isolationoxide] (0,0.75) rectangle (1,2);
\fill[isolationoxide] (8.5,0.75) rectangle (11.5,2);
\fill[isolationoxide] (19,0.75) rectangle (20,2);
% n-well
\fill[nwell] (1,0.75) rectangle (8.5,2);
\node at (5.75,1) {N-Well};

% gate oxide
\fill[LightGray] (4.8,2) rectangle (6.7,2.3);
\fill[LightGray] (13.3,2) rectangle (15.2,2.3);
% gate poly
\fill[BrickRed] (4.8,2.3) rectangle (6.7,2.6);
\fill[BrickRed] (13.3,2.3) rectangle (15.2,2.6);
	\end{tikzpicture}
	\begin{tikzpicture}[node distance = 3cm, auto, thick,scale=\CrossAndTopSection, every node/.style={transform shape}]
		\fill[YellowOrange] (0,0) rectangle (20,10);
% n-well
\fill[Goldenrod] (1.25,1.5) rectangle (8.25,7.25);
% field oxide
\fill[DarkGray] (0,0) rectangle (1,12);
\fill[DarkGray] (8.5,0) rectangle (11.5,12);
\fill[DarkGray] (19,0) rectangle (20,12);
\fill[DarkGray] (0,0) rectangle (20,1.25);
\fill[DarkGray] (0,7.5) rectangle (20,12);
% poly
\fill[BrickRed] (4.8,1.75) rectangle (6.7,9);
\fill[BrickRed] (13.3,1.75) rectangle (15.2,9);
\fill[BrickRed] (4.8,8) rectangle (15.2,9);
	\end{tikzpicture}
	\caption{Resist removal}
\end{figure}